% ============================================
% 四、程序完善题 — 讲稿
% ============================================
\section{程序完善题}

程序完善题是试卷中第四大题，共3小题、10个空、每空2分，合计20分。
与阅读程序填空题不同，完善题要求你\underline{根据上下文逻辑自己填写缺失的语句}，
考查的是对C语言基本语法模式的\underline{熟练度}和\underline{模式识别能力}。

\textbf{教学提示：} 强调完善题的策略 —— 不要想太多，每道题对应的是很直白的语法点。先通读题目理解意图，然后观察空白周围代码的上下文，判断空白处需要完成什么操作。

\bigskip

% ====================================
\subsection{第1题：字符类型判断（3空）}

\textbf{考查知识点：} getchar() 函数、字符ASCII码范围判断、if-else if 多分支结构。

\subsubsection{题目原文}

从标准输入读入一个字符，判断该字符是字母、数字、空格还是其他字符，并输出相应提示。

\subsubsection{代码还原}

\begin{lstlisting}[language=C]
char ch;
ch = 【1】;                                   // 读入字符
if (【2】)                                    // 判断是否为字母
    printf("It is an English character\n");
else if (【3】)                               // 判断是否为数字
    printf("It is a digit character\n");
else if (ch == ' ')
    printf("It is a space character\n");
else
    printf("It is other character\n");
\end{lstlisting}

\subsubsection{逐空讲解}

\textbf{【1】\texttt{getchar()}} —— 从标准输入读取一个字符。

\vspace{0.3em}
\textbf{这个空考什么：} 字符输入函数的调用方式。很多同学看到输入就想写 \texttt{scanf("\%c", \&ch)}，这在语法上也可以，但题目上下文只有一行调用、且返回结果直接赋值给 \texttt{ch}，明显是 \texttt{getchar()} 的用法模式。scanf 需要格式化串和取地址，不适合这个位置。

\textbf{常见错误：}
\begin{itemize}
    \item \texttt{scanf("\%c", \&ch)} —— 没有错，但这空只需填一行，scanf在此处显得多余。关键是题目就是想考你 \texttt{getchar()} 这个函数。
    \item \texttt{getch()} —— 不是标准C函数，是 conio.h 里的非标准函数，考试不认。
    \item \texttt{getche()} —— 同上，非标准。
\end{itemize}

\vspace{0.6em}
\textbf{【2】\texttt{(ch>='a' \&\& ch<='z') || (ch>='A' \&\& ch<='Z')}} —— 判断字符是否为英文字母。

\vspace{0.3em}
\textbf{这个空考什么：} 字符的ASCII码范围比较。C语言中字符本质上是一个整数（ASCII码值），所以可以直接用 \texttt{>=}、\texttt{<=} 比较。

\textbf{分两步理解：}
\begin{enumerate}
    \item \texttt{ch>='a' \&\& ch<='z'} —— ch 的 ASCII 码在小写字母 a(97) 到 z(122) 之间。
    \item \texttt{ch>='A' \&\& ch<='Z'} —— ch 的 ASCII 码在大写字母 A(65) 到 Z(90) 之间。
    \item 用 \texttt{||} 连接 —— 因为字母包含大写和小写，两者是「或」的关系。
\end{enumerate}

\textbf{常见错误：}
\begin{itemize}
    \item 写成 \texttt{'a'<=ch<='z'} —— \textbf{这是经典错误！}C语言不支持数学中的连续比较。\texttt{'a'<=ch<='z'} 实际上等价于 \texttt{('a'<=ch)<='z'}，先算 \texttt{'a'<=ch} 得到 0 或 1，再和 \texttt{'z'} 比较，结果永远是 1，逻辑完全错误。
    \item 漏掉大写范围 —— 只写 \texttt{ch>='a' \&\& ch<='z'}，会漏掉大写字母。
    \item 把 \texttt{||} 写成 \texttt{\&\&} —— 大写和小写不可能同时满足，用 \texttt{\&\&} 条件永远为假。
    \item 用 \texttt{isalpha(ch)} —— 函数调用确实可以，但需要 \texttt{\#include <ctype.h>}。完善题一般考的是你手写判断的能力。
\end{itemize}

\vspace{0.6em}
\textbf{【3】\texttt{ch>='0' \&\& ch<='9'}} —— 判断字符是否为数字字符。

\vspace{0.3em}
\textbf{这个空考什么：} 字符数字的ASCII码范围。注意「数字字符」和「数值」的区别：
\begin{itemize}
    \item \texttt{ch>='0' \&\& ch<='9'} —— 判断 ch 是字符 \texttt{'0'} 到 \texttt{'9'}
    \item 不是判断 ch 本身的值在 0 到 9 之间！
    \item 字符 \texttt{'0'} 的ASCII码是48，不是0。
\end{itemize}

\textbf{常见错误：}
\begin{itemize}
    \item \texttt{ch>=0 \&\& ch<=9} —— 这是在比较ASCII码值 0\~{}9（控制字符），不是数字字符。
    \item \texttt{ch>=48 \&\& ch<=57} —— 语法上没错（这是数字字符的ASCII范围），但阅卷老师会认为你没有理解字符比较的本质。写 \texttt{'0'} 和 \texttt{'9'} 更清晰、更可读。
    \item 忘记引号，写成 \texttt{ch>=0 \&\& ch<=9} —— 同上。
\end{itemize}

\subsubsection{教学笔记}

这道题的三个空分别考查了 \texttt{getchar()}、字母判断、数字判断，是C语言字符处理的基础三件套。
考试中这类题频率很高，建议记住：
\begin{itemize}
    \item 字母：\texttt{(ch>='a' \&\& ch<='z') || (ch>='A' \&\& ch<='Z')}
    \item 数字：\texttt{ch>='0' \&\& ch<='9'}
    \item 小写字母：\texttt{ch>='a' \&\& ch<='z'}
    \item 大写字母：\texttt{ch>='A' \&\& ch<='Z'}
    \item 大小写转换：\texttt{ch - 'a' + 'A'}（小写转大写）
\end{itemize}

\textbf{教学提示：} 可以当场问学生：「\texttt{'5'} 减去 \texttt{'0'} 等于多少？」答案应该是 5 —— 字符 \texttt{'5'} 的 ASCII 码 53 减去 \texttt{'0'} 的 48 等于 5。这是字符转数字的常用技巧。

\bigskip
\bigskip

% ====================================
\subsection{第2题：数组查找（2空）}

\textbf{考查知识点：} scanf 地址传递、数组元素比较、break 结束循环。

\subsubsection{题目原文}

在一个包含10个整数的数组中，查找与给定值 x 相同的元素所在位置（下标）。

\subsubsection{代码还原}

\begin{lstlisting}[language=C]
int a[10], i, x;
printf("input 10 integers:");
for (i = 0; i < 10; i++) scanf("%d", 【4】);   // 输入数组元素
printf("input the number you want to find x:");
scanf("%d", &x);
for (i = 0; i < 10; i++)
    if (【5】) break;                           // 找到则跳出循环
if (i < 10)
    printf("the pos of x is: %d\n", i);
else
    printf("can not find x!\n");
\end{lstlisting}

\subsubsection{逐空讲解}

\textbf{【4】\texttt{\&a[i]}} —— scanf 需要变量的地址。

\vspace{0.3em}
\textbf{这个空考什么：} scanf 函数的取地址语法。这是初学者最易出错的知识点之一。

\textbf{核心概念：}
\begin{itemize}
    \item scanf 的作用是把键盘输入的值写入变量，所以它需要知道变量在哪里——即变量的\underline{内存地址}。
    \item \texttt{\&} 是取地址运算符，\texttt{\&a[i]} 表示数组第 i 个元素的地址。
    \item 等价写法：\texttt{a+i}（数组名 a 是首地址，+i 偏移 i 个元素）。
\end{itemize}

\textbf{常见错误：}
\begin{itemize}
    \item \texttt{a[i]} —— \textbf{最高频错误！} scanf 传的是值而不是地址。对于整型变量，scanf("\%d", a[i]) 会把 a[i] 的\underline{未初始化值}当作地址传给 scanf，导致段错误或写入非法内存。
    \item \texttt{a} —— 每次循环都给 scanf 传同一个地址 a（即 a[0] 的地址），结果所有输入都覆盖在 a[0] 上。虽然 \texttt{\&a[0]} 和 \texttt{a} 地址相同，但写成 a 就失去了循环的意义。
    \item \texttt{\&a} —— 取整个数组的地址（类型是 int(*)[10]），类型不对，虽然地址值可能相同但语义错误。
\end{itemize}

\textbf{记忆技巧：} scanf 的参数永远是\underline{地址}。数组元素用 \texttt{\&a[i]}，单个变量用 \texttt{\&x}，字符串用 \texttt{str}（数组名本身就是地址）。

\vspace{0.6em}
\textbf{【5】\texttt{a[i]==x}} —— 判断当前数组元素是否等于目标值。

\vspace{0.3em}
\textbf{这个空考什么：} 相等比较运算符。C 语言中 \texttt{==} 是比较，\texttt{=} 是赋值。

\textbf{为什么这里用 break：}
\begin{itemize}
    \item 一旦找到目标值，i 就是目标的下标，不需要继续循环。
    \item break 直接跳出 for 循环，i 保留找到时的值。
    \item 后面的 \texttt{if (i < 10)} 判断就是利用了这个特性：如果正常找到了，i 一定 < 10；如果循环跑完也没 break，i 会是10。
\end{itemize}

\textbf{常见错误：}
\begin{itemize}
    \item \texttt{a[i]=x} —— \textbf{致命错误！} 把 \texttt{x} 的值赋给了 \texttt{a[i]}，同时赋值表达式的结果是 x 的值（非零即真），if 条件永远为真，第一次循环就会 break。修改了原数组数据，且逻辑完全错误。
    \item \texttt{a[i]==x \&\& break} —— break 不能放在表达式中，语法错误。
\end{itemize}

\subsubsection{教学笔记}

这道题的核心是区分 \texttt{\&a[i]} 和 \texttt{a[i]} 在使用场景上的差异：
\begin{itemize}
    \item scanf/\textbf{输入}时 —— 需要 \texttt{\&a[i]}（地址）
    \item 计算/判断/\textbf{输出}时 —— 需要 \texttt{a[i]}（值）
\end{itemize}

\textbf{教学提示：} 可以在黑板上画一个内存图，展示 a[0]\~{}a[9] 各占4字节，\texttt{\&a[0]} 指向第0个元素的起始地址。让学生直观理解「地址」和「值」的区别。

\bigskip
\bigskip

% ====================================
\subsection{第3题：统计非负数（5空）}

\textbf{考查知识点：} 变量初始化、continue 跳过循环、累加器模式、printf 格式化输出。

\subsubsection{题目原文}

读入20个整数，统计其中非负数的个数，并计算非负数之和。

\subsubsection{代码还原}

\begin{lstlisting}[language=C]
int i, a[20], s, count;
s = 【6】;                              // 累加和初始化
for (i = 0; i < 20; i++) scanf("%d", &a[i]);
for (i = 0; i < 20; i++) {
    if (【7】) continue;                // 跳过负数
    【8】                                // 累加非负数
    【9】                                // 计数加1
}
printf("s = %d\t count = %d\n", 【10】); // 输出结果
\end{lstlisting}

\subsubsection{逐空讲解}

\textbf{【6】\texttt{s=0}} —— 累加和变量初始化。

\vspace{0.3em}
\textbf{这个空考什么：} C 语言局部变量\underline{不会自动初始化为0}。如果不手动赋值，s 的值是内存中的垃圾数据，累加结果必然错误。

\textbf{常见错误：}
\begin{itemize}
    \item 不填 —— 认为 s 默认就是 0，这是 \textbf{Java/Python 思维}，不适用于 C。
    \item 写 \texttt{count=0} —— count 初始化也是需要的，但根据上下文，【6】的位置紧接 s，且后续代码中 count 在循环内才递增，所以【6】处先初始化 s。不过题目实际是给【6】处的 s 一个初始值。
\end{itemize}

\textbf{注意：} 观察第2行声明 \texttt{int i, a[20], s, count;} 之后 s 和 count 都没有被赋予初始值。在 C 语言中，非静态局部变量的初值是\underline{不确定}的。所以必须在使用前初始化。s 和 count 都需要初始化，本题中 s 在循环外初始化为0，count 在循环内通过 count++ 计数（但 count 本身也需要初始化为0，通常在声明后、循环前完成。这里的题目排版略有不严谨，但习惯上 s 和 count 都应初始化为0）。实际考试中注意观察代码上下文来确定是 s=0 还是 count=0。

\vspace{0.6em}
\textbf{【7】\texttt{a[i]<0}} —— 判断当前元素是否为负数。

\vspace{0.3em}
\textbf{这个空考什么：} continue 语句的用法 + 负数条件判断。

\textbf{逻辑分析：}
\begin{itemize}
    \item 题目要求统计\underline{非负数}的个数和求和。
    \item 所以遇到负数应该\underline{跳过}，不做任何处理。
    \item continue 的作用是结束本次循环迭代，直接进入 i++ 开始下一轮。
    \item 因此 if 的条件应该是「是负数」→ \texttt{a[i] < 0}。
\end{itemize}

\textbf{常见错误：}
\begin{itemize}
    \item \texttt{a[i]>=0} —— 逻辑反了！如果条件写成非负数时才 continue，那么非负数被跳过、负数被统计，结果完全相反。
    \item \texttt{a[i]<=0} —— 把 0 也跳过了，但 0 是非负数，应该统计。
\end{itemize}

\vspace{0.6em}
\textbf{【8】\texttt{s+=a[i];} / \texttt{s=s+a[i];}} —— 累加非负数到总和。

\vspace{0.3em}
\textbf{这个空考什么：} 复合赋值运算符 \texttt{+=}（或者普通赋值加法）。

\textbf{位置分析：} continue 之后才执行到这里，意味着执行到【8】的元素一定是非负数。所以这里直接累加即可。

\textbf{常见错误：}
\begin{itemize}
    \item \texttt{s = a[i]} —— 覆盖而非累加，最终 s 只会保存最后一个非负数的值。
    \item \texttt{s + a[i];} —— 表达式有值但没有赋值，s 不会改变。
\end{itemize}

\vspace{0.6em}
\textbf{【9】\texttt{count++;} / \texttt{++count;} / \texttt{count+=1;} / \texttt{count=count+1;}} —— 非负数计数加一。

\vspace{0.3em}
\textbf{这个空考什么：} 自增运算符 \texttt{++} 的用法。

\textbf{注意【8】和【9】的顺序：} 可以先累加再计数，也可以先计数再累加，不影响结果。考试中两空都给分（只要逻辑正确）。

\textbf{常见错误：}
\begin{itemize}
    \item \texttt{count+1;} —— 没有赋值，count 不变。
    \item \texttt{count = 1;} —— 每次循环都把 count 设为 1，最终 count 永远是 1。
\end{itemize}

\vspace{0.6em}
\textbf{【10】\texttt{s, count}} —— printf 的参数列表。

\vspace{0.3em}
\textbf{这个空考什么：} printf 格式化字符串与参数的一一对应关系。

格式化串中有两个占位符：\texttt{\%d} 和 \texttt{\%d}，分别对应格式串后的第1个和第2个参数。所以这里应填 \texttt{s, count}（注意顺序：先 s 后 count，与格式串中 \texttt{\%d} 出现的顺序一致）。

\textbf{常见错误：}
\begin{itemize}
    \item \texttt{count, s} —— 格式串中 "s =" 在前、"count =" 在后，如果参数顺序颠倒，输出的数字会张冠李戴。
    \item \texttt{\&s, \&count} —— printf 不需要取地址！\textbf{scanf 要地址，printf 要值。}这是学生最容易混淆的一组知识点。
\end{itemize}

\subsubsection{教学笔记}

这道题的5个空涵盖了变量初始化、循环控制、累加器模式、自增运算、格式化输出，构成了一个完整的「遍历数组 + 条件筛选 + 统计」范式。这个范式在 C 语言编程中非常常见，建议学生熟练掌握。

\textbf{教学提示：} 可以带领学生把代码逻辑用自然语言复述一遍：「初始化 s 为 0 → 读入20个数 → 遍历数组 → 如果当前数是负数就跳过（continue）→ 否则把它加到 s 上，count 加一 → 最后输出 s 和 count。」让学生养成「读代码→翻译成自然语言→理解逻辑」的习惯。

\bigskip
\bigskip

% ====================================
\subsection{程序完善题总结}

\subsubsection{题型特点}

程序完善题与阅读填空题的最大区别：
\begin{itemize}
    \item 阅读填空题：代码是完整的，你在分析已有代码的运行结果。
    \item 完善题：代码有空白，你需要自己写出缺失的语句。
\end{itemize}

完善题考查的不是你「能不能看懂了」，而是你「会不会写」。

\subsubsection{答题策略}

\begin{enumerate}
    \item \textbf{先通读整段代码} —— 理解程序要完成什么功能。
    \item \textbf{观察空白周围代码} —— 从上下文推断空白处需要什么类型的语句。
    \item \textbf{确定空白的基本语法模式} —— 是函数调用（getchar）？是条件判断（if）？是赋值（=）？是输出（printf的参数）？
    \item \textbf{填写后检查} —— 把填好的代码在脑中运行一遍，验证是否合理。
\end{enumerate}

\subsubsection{本大题考点汇总}

\begin{table}[h]
    \centering
    \begin{tabular}{@{}cll@{}}
        \toprule
        \textbf{空号} & \textbf{答案}                          & \textbf{知识点}        \\
        \midrule
        【1】         & \texttt{getchar()}                     & 字符输入函数              \\
        【2】         & \texttt{(ch>='a'\&\&ch<='z')||...}    & ASCII 范围判断          \\
        【3】         & \texttt{ch>='0' \&\& ch<='9'}         & 数字字符判断              \\
        【4】         & \texttt{\&a[i]}                        & scanf 地址传递          \\
        【5】         & \texttt{a[i]==x}                       & 相等比较（区分 = 和 ==）     \\
        【6】         & \texttt{s=0}                           & 变量初始化               \\
        【7】         & \texttt{a[i]<0}                        & continue 跳过条件       \\
        【8】         & \texttt{s+=a[i];}                      & 复合赋值 / 累加器模式        \\
        【9】         & \texttt{count++;}                      & 自增运算符 / 计数模式        \\
        【10】        & \texttt{s, count}                      & printf 参数对应          \\
        \bottomrule
    \end{tabular}
\end{table}

\textbf{教学提示：} 这个表可以在讲完后投影出来，让学生快速回顾所有考点。也可以提问学生：「哪个空你最容易填错？」根据反馈做针对性强调。
